The framework's thesis is that binary validity is not enough: two models can
both fail and still fail for different reasons. The metric set therefore splits
planning quality into retry dependence, executable structure, schema grounding,
precondition awareness, and reasoning-plan alignment. Each dimension is tied to
a specific failure mechanism reported in the LLM-planning literature.

\begin{quote}\small
If LLM planning is closer to approximate retrieval than to search, retries
should not be treated as free evidence of planning ability. First-attempt
success, retry efficiency, executable prefixes, hallucination rates, and
reasoning-faithfulness must be read together.
\end{quote}

\begin{table}[h]\small
\centering
\begin{tabularx}{\linewidth}{@{}p{3.2cm} p{3.1cm} >{\raggedright\arraybackslash}X@{}}
\toprule
\textbf{Literature claim} & \textbf{Metric signal} & \textbf{Design consequence}\\
\midrule
LLMs may solve planning tasks by retrieving surface patterns rather than by
performing explicit search. & \textbf{FASR}, \textbf{Retry Gap},
\textbf{CoT Alignment} & First-attempt success is separated from final success,
and reasoning traces are checked against extracted plan actions instead of being
accepted as explanations.\\
\addlinespace
Feedback and repeated attempts can behave like stochastic re-sampling rather
than correction. & \textbf{IWSR}, \textbf{Retry Gap},
\textbf{$P(\mathrm{Valid}\mid k)$ profile} & A model that succeeds only after
many attempts is discounted; early repairs count more than late lucky samples.\\
\addlinespace
Grounding can fail before sequencing begins: the model can invent action names
or objects not present in the PDDL task. & \textbf{Hallucination rate},
\textbf{Fuzzy hallucination rate}, \textbf{Object hallucination rate},
\textbf{IHR} & Action vocabulary and object grounding are measured separately
from VAL validity, so invalid plans can still be diagnosed precisely.\\
\addlinespace
Some failures are local state-tracking errors rather than total domain
ignorance. & \textbf{Executability ratio}, \textbf{PAS},
\textbf{Mean temporal distance} & The simulator measures how far the plan can
run and classifies the first failed precondition as sequencing error or state
fabrication.\\
\addlinespace
Capability profiles need multiple signals, because the same success rate can
hide different mechanisms. & \textbf{SR}, \textbf{FASR}, \textbf{IWSR},
\textbf{RG}, \textbf{IHR}, \textbf{ER}, \textbf{PAS}, \textbf{CoT} & The report
classifies models as genuine planners, efficient correctors, stochastic
searchers, lucky retrievers, understanders, vocabulary-only systems, or
ungrounded systems.\\
\bottomrule
\end{tabularx}
\caption{Scientific motivation for the metric families used in Section~7.}
\label{tab:section07_metric_foundation}
\end{table}

\paragraph{Retry-aware scoring.}
Success rate alone is intentionally not the primary capability estimate. The
analysis computes \textbf{FASR} from the private \code{_iter1} column and
\textbf{IWSR} from \code{_iwsr_contrib}. This keeps three cases separate:
models that solve immediately, models that repair quickly, and models whose
reported success is mostly a product of repeated sampling.

\paragraph{Grounding before validity.}
The hallucination metrics measure whether the generated plan is expressed in the
task's legal vocabulary. Strict action hallucination catches invented operators;
fuzzy hallucination avoids over-penalising small naming variants; object
hallucination checks whether arguments come from the problem's declared objects.
This follows the idea that a model cannot be credited with planning over a
formal domain if it is not grounded in that domain's symbols.

\paragraph{Execution as a diagnostic, not only a pass/fail gate.}
The pure-Python \code{PDDLSimulator} replays the representative plan and records
the prefix that remains executable. At the first failed precondition, a
previously true atom is counted as a sequencing error; an atom that was never
true is counted as state fabrication. \textbf{PAS} and temporal distance
therefore distinguish near-miss ordering failures from invented world states.

\paragraph{Reasoning-faithfulness.}
CoTaL is used only when reasoning traces are available. It compares actions
mentioned in the reasoning channel with actions extracted from the final plan,
so chain-of-thought is treated as evidence only when it is aligned with the
actual output. This prevents fluent but post-hoc rationales from inflating the
planning interpretation.

\paragraph{Composite score.}
The composite Planning Score keeps the same bias as the metric design:
\[
PS = 0.25 \cdot FASR
   + 0.20 \cdot IWSR
   + 0.20 \cdot ER
   + 0.20 \cdot (1 - Hallucination)
   + 0.15 \cdot PAS .
\]
FASR has the largest weight because it cannot be inflated by retries. IWSR and
ER capture repair efficiency and executable structure; grounding is represented
by inverse hallucination; PAS receives a smaller weight because it is undefined
for perfectly executable plans and requires a neutral prior when aggregated.

\begin{table}[h]\small
\centering
\begin{tabularx}{\linewidth}{@{}p{3.0cm} p{4.2cm} >{\raggedright\arraybackslash}X@{}}
\toprule
\textbf{Reference family} & \textbf{Used for} & \textbf{Where it appears in the framework}\\
\midrule
Kambhampati (2024) & Retrieval-vs-planning framing; surface-pattern solutions;
post-hoc reasoning caveats. & Hallucination rationale, lucky-retriever profile,
CoT interpretation.\\
Valmeekam et al. (2023) & Formal-domain grounding tests and obfuscation-style
checks. & IHR thresholds, no-grounding profile, understander profile.\\
Stechly et al. (2023) & Retry behaviour and repeated sampling under feedback. &
Retry Gap, IWSR, stochastic-searcher profile, $P(\mathrm{Valid}\mid k)$.\\
Dziri et al. (2023) & Compositional and long-horizon state-tracking failures. &
PAS, temporal distance, understander profile.\\
Vafa et al. (2024) and McCoy et al. (2024) & World-model incoherence and
training-frequency / distributional explanations. & Vocabulary-only and
no-grounding interpretations.\\
\bottomrule
\end{tabularx}
\caption{Reference families used by the metric taxonomy.}
\label{tab:section07_reference_families}
\end{table}
